\textbf{Phần 1. Trắc nghiệm lựa chọn}
\setcounterpageref{ex}{0}

\begin{ex} %[3H1V19]
Hình vuông có đặc điểm nổi bật nào sau đây?
\choice
{\True Có 4 góc vuông và 4 cạnh có độ dài bằng nhau}
{Có 4 góc vuông và 2 cạnh dài bằng nhau}
{Có 3 góc vuông và 4 cạnh bằng nhau}
{Có 4 cạnh bằng nhau nhưng không có góc vuông}
\loigiai{
Hình vuông có 4 góc vuông và 4 cạnh dài bằng nhau.
\textbf{Đáp án của câu hỏi là A}
}
\end{ex}

\begin{ex} %[3H1V19]
Cho một hình tam giác có độ dài các cạnh lần lượt là $15\text{ cm}, 20\text{ cm}, 25\text{ cm}$. Chu vi của hình tam giác đó là
\choice
{$50\text{ cm}$}
{$55\text{ cm}$}
{\True $60\text{ cm}$}
{$65\text{ cm}$}
\loigiai{
Chu vi hình tam giác là: $15 + 20 + 25 = 60\text{ (cm)}$.
\textbf{Đáp án của câu hỏi là C}
}
\end{ex}

\begin{ex} %[3H1V19]
Một hình chữ nhật có chiều dài $24\text{ cm}$, chiều rộng bằng một nửa chiều dài. Chu vi hình chữ nhật đó là
\choice
{$48\text{ cm}$}
{\True $72\text{ cm}$}
{$36\text{ cm}$}
{$60\text{ cm}$}
\loigiai{
Chiều rộng hình chữ nhật: $24 : 2 = 12\text{ cm}$. Chu vi: $(24 + 12) \times 2 = 72\text{ cm}$.
\textbf{Đáp án của câu hỏi là B}
}
\end{ex}

\begin{ex} %[3H1V19]
Một sợi dây thép dài $36\text{ cm}$ được uốn thành một hình vuông. Độ dài mỗi cạnh của hình vuông đó là
\choice
{$6\text{ cm}$}
{$12\text{ cm}$}
{$8\text{ cm}$}
{\True $9\text{ cm}$}
\loigiai{
Độ dài mỗi cạnh của hình vuông là: $36 : 4 = 9\text{ (cm)}$.
\textbf{Đáp án của câu hỏi là D}
}
\end{ex}

\begin{ex} %[3H1V19]
Cho một hình chữ nhật có chu vi bằng chu vi của hình vuông cạnh $8\text{ cm}$. Biết chiều rộng của hình chữ nhật là $6\text{ cm}$. Chiều dài của hình chữ nhật đó là
\choice
{\True $10\text{ cm}$}
{$12\text{ cm}$}
{$8\text{ cm}$}
{$14\text{ cm}$}
\loigiai{
Chu vi hình vuông là: $8 \times 4 = 32\text{ (cm)}$. Nửa chu vi hình chữ nhật là: $32 : 2 = 16\text{ (cm)}$. Chiều dài hình chữ nhật là: $16 - 6 = 10\text{ (cm)}$.
\textbf{Đáp án của câu hỏi là A}
}
\end{ex}

\begin{ex} %[3H1V19]
Người ta ghép hai hình vuông nhỏ cạnh $5\text{ cm}$ cạnh nhau để tạo thành một hình chữ nhật. Chu vi của hình chữ nhật tạo thành là
\choice
{$40\text{ cm}$}
{$25\text{ cm}$}
{\True $30\text{ cm}$}
{$35\text{ cm}$}
\loigiai{
Hình chữ nhật có chiều dài $5 + 5 = 10\text{ cm}$ và chiều rộng $5\text{ cm}$. Chu vi là: $(10 + 5) \times 2 = 30\text{ (cm)}$.
\textbf{Đáp án của câu hỏi là C}
}
\end{ex}

\begin{ex} %[3H1V19]
Khẳng định nào sau đây là ĐÚNG về mối quan hệ giữa hình vuông và hình chữ nhật?
\choice
{Mỗi hình chữ nhật đều là một hình vuông}
{\True Mỗi hình vuông đều là một hình chữ nhật đặc biệt có 4 cạnh bằng nhau}
{Hình chữ nhật có 4 cạnh bằng nhau}
{Hình vuông không có góc vuông}
\loigiai{
Hình vuông có đủ 4 góc vuông và có các cặp cạnh đối diện bằng nhau (thậm chí cả 4 cạnh đều bằng nhau) nên hình vuông là trường hợp đặc biệt của hình chữ nhật.
\textbf{Đáp án của câu hỏi là B}
}
\end{ex}

\begin{ex} %[3H1V19]
Cho tam giác $ABC$ có cạnh $AB = AC = 18\text{ cm}$, cạnh $BC$ ngắn hơn cạnh $AB$ là $6\text{ cm}$. Chu vi của tam giác $ABC$ là
\choice
{$42\text{ cm}$}
{$54\text{ cm}$}
{$36\text{ cm}$}
{\True $48\text{ cm}$}
\loigiai{
Cạnh $BC$ dài: $18 - 6 = 12\text{ (cm)}$. Chu vi tam giác $ABC$ là: $18 + 18 + 12 = 48\text{ (cm)}$.
\textbf{Đáp án của câu hỏi là D}
}
\end{ex}

\begin{ex} %[3H1V19]
Một mảnh vườn hình chữ nhật có chiều dài $30\text{ m}$, nếu giảm chiều dài đi $6\text{ m}$ thì mảnh vườn trở thành hình vuông. Chu vi của mảnh vườn ban đầu là
\choice
{\True $108\text{ m}$}
{$96\text{ m}$}
{$120\text{ m}$}
{$84\text{ m}$}
\loigiai{
Khi giảm chiều dài 6 m thành hình vuông thì cạnh hình vuông là: $30 - 6 = 24\text{ m}$. Chiều rộng mảnh vườn chính bằng $24\text{ m}$. Chu vi ban đầu là: $(30 + 24) \times 2 = 108\text{ (m)}$.
\textbf{Đáp án của câu hỏi là A}
}
\end{ex}

\begin{ex} %[3H1V19]
Một hình tứ giác có độ dài 4 cạnh là các số tự nhiên liên tiếp tăng dần. Biết chu vi của hình tứ giác đó là $34\text{ cm}$. Cạnh dài nhất của hình tứ giác là
\choice
{$9\text{ cm}$}
{$11\text{ cm}$}
{\True $10\text{ cm}$}
{$8\text{ cm}$}
\loigiai{
Gọi 4 cạnh là $a, a+1, a+2, a+3$. Tổng chu vi: $4 \times a + 6 = 34 \Rightarrow 4 \times a = 28 \Rightarrow a = 7$. Cạnh dài nhất là $7 + 3 = 10\text{ cm}$.
\textbf{Đáp án của câu hỏi là C}
}
\end{ex}

\begin{ex} %[3H1V19]
Có bao nhiêu đoạn thẳng nối từ đỉnh một tứ giác đến các đỉnh còn lại không kể các cạnh của tứ giác đó?
\choice
{$2$ đoạn thẳng}
{\True $1$ đoạn thẳng (đường chéo)}
{$3$ đoạn thẳng}
{$4$ đoạn thẳng}
\loigiai{
Từ 1 đỉnh của tứ giác, nối đến 2 đỉnh liền kề là 2 cạnh, chỉ có 1 đỉnh đối diện tạo thành 1 đường chéo.
\textbf{Đáp án của câu hỏi là B}
}
\end{ex}

\begin{ex} %[3H1V19]
Người ta cắt một hình chữ nhật có chiều dài $15\text{ cm}$, chiều rộng $9\text{ cm}$ để lấy ra một hình vuông có cạnh lớn nhất có thể. Chu vi của phần hình chữ nhật còn lại là
\choice
{$24\text{ cm}$}
{$36\text{ cm}$}
{$18\text{ cm}$}
{\True $30\text{ cm}$}
\loigiai{
Hình vuông lớn nhất có cạnh bằng chiều rộng là $9\text{ cm}$. Phần còn lại là hình chữ nhật có kích thước $9\text{ cm}$ và $15 - 9 = 6\text{ cm}$. Chu vi phần còn lại là: $(9 + 6) \times 2 = 30\text{ (cm)}$.
\textbf{Đáp án của câu hỏi là D}
}
\end{ex}

\textbf{Phần 2. Trắc nghiệm đúng sai}
\setcounterpageref{ex}{0}

\begin{ex} %[3H1V19]
Xét tính đúng sai của các nhận định về các hình phẳng:
\choiceTF
{\True Hình vuông có 4 góc vuông và 4 cạnh dài bằng nhau.}
{Mọi hình tứ giác đều có 4 góc vuông.}
{\True Tam giác có đúng 3 cạnh và 3 góc.}
{Hình chữ nhật có 4 cạnh bằng nhau.}
\loigiai{
\begin{itemize}
    \item Ý a) Đúng: Đặc điểm định nghĩa của hình vuông.
    \item Ý b) Sai: Chỉ có hình vuông và hình chữ nhật mới có 4 góc vuông.
    \item Ý c) Đúng: Định nghĩa hình tam giác.
    \item Ý d) Sai: Hình chữ nhật nói chung chỉ có các cặp cạnh đối diện bằng nhau.
\end{itemize}
\begin{itemize}
    \item a) \textbf{Đúng}.
    \item b) \textbf{Sai}.
    \item c) \textbf{Đúng}.
    \item d) \textbf{Sai}.
\end{itemize}
}
\end{ex}

\begin{ex} %[3H1V19]
Một hình chữ nhật có chiều dài $16\text{ cm}$, chiều rộng $8\text{ cm}$:
\choiceTF
{\True Chiều dài gấp đôi chiều rộng.}
{\True Chu vi của hình chữ nhật đó là $48\text{ cm}$.}
{Nếu gấp đôi chiều rộng thì hình chữ nhật trở thành hình vuông có cạnh $8\text{ cm}$.}
{\True Có thể chia hình chữ nhật đó thành 2 hình vuông cạnh $8\text{ cm}$.}
\loigiai{
\begin{itemize}
    \item Ý a) Đúng: $16 : 8 = 2$ lần.
    \item Ý b) Đúng: $(16 + 8) \times 2 = 48\text{ cm}$.
    \item Ý c) Sai: Nếu gấp đôi chiều rộng thì chiều rộng thành $16\text{ cm}$, bằng chiều dài nên thành hình vuông cạnh $16\text{ cm}$, không phải $8\text{ cm}$.
    \item Ý d) Đúng: $16 : 8 = 2$ hình vuông.
\end{itemize}
\begin{itemize}
    \item a) \textbf{Đúng}.
    \item b) \textbf{Đúng}.
    \item c) \textbf{Sai}.
    \item d) \textbf{Đúng}.
\end{itemize}
}
\end{ex}

\textbf{Phần 3. Trả lời ngắn}
\setcounterpageref{ex}{0}

\begin{ex} %[3H1V19]
Tính chu vi hình tam giác đều có độ dài mỗi cạnh bằng $15\text{ cm}$.
\shortans{45}
\loigiai{
\textbf{Đáp số:} $45$\\
Chu vi là: $15 \times 3 = 45\text{ (cm)}$.
}
\end{ex}

\begin{ex} %[3H1V19]
Một hình vuông có chu vi $64\text{ cm}$. Tính độ dài cạnh của hình vuông đó theo xăng-ti-mét.
\shortans{16}
\loigiai{
\textbf{Đáp số:} $16$\\
Cạnh hình vuông là: $64 : 4 = 16\text{ (cm)}$.
}
\end{ex}

\begin{ex} %[3H1V19]
Một hình chữ nhật có chu vi $50\text{ cm}$, chiều dài là $18\text{ cm}$. Tính chiều rộng của hình chữ nhật theo xăng-ti-mét.
\shortans{7}
\loigiai{
\textbf{Đáp số:} $7$\\
Nửa chu vi: $50 : 2 = 25\text{ cm}$. Chiều rộng: $25 - 18 = 7\text{ (cm)}$.
}
\end{ex}

\begin{ex} %[3H1V19]
Một mảnh đất hình chữ nhật có chiều dài $25\text{ m}$, chiều rộng ngắn hơn chiều dài $7\text{ m}$. Tính chu vi của mảnh đất đó theo mét.
\shortans{86}
\loigiai{
\textbf{Đáp số:} $86$\\
Chiều rộng: $25 - 7 = 18\text{ m}$. Chu vi: $(25 + 18) \times 2 = 86\text{ (m)}$.
}
\end{ex}

\begin{ex} %[3H1V19]
Bốn que tính dài bằng nhau ghép thành một hình vuông có chu vi $32\text{ cm}$. Hỏi dùng 6 que tính như vậy ghép thành một hình chữ nhật thì chu vi của hình chữ nhật ghép được là bao nhiêu xăng-ti-mét?
\shortans{48}
\loigiai{
\textbf{Đáp số:} $48$\\
Mỗi que tính dài: $32 : 4 = 8\text{ cm}$. Chu vi hình ghép bởi 6 que tính là: $6 \times 8 = 48\text{ (cm)}$.
}
\end{ex}

\begin{ex} %[3H1V19]
Người ta ghép 3 hình vuông bằng nhau cạnh $4\text{ cm}$ thành một hàng ngang để được một hình chữ nhật. Tính chu vi hình chữ nhật đó theo xăng-ti-mét.
\shortans{32}
\loigiai{
\textbf{Đáp số:} $32$\\
Hình chữ nhật có chiều dài $4 \times 3 = 12\text{ cm}$, chiều rộng $4\text{ cm}$. Chu vi: $(12 + 4) \times 2 = 32\text{ (cm)}$.
}
\end{ex}
